% Generated table; it uses the manuscript preamble's macros (\hltint, \hcell, \altrow, PanelBg).
\begin{table}[pos=htbp]
\tableserif
\centering
\begin{threeparttable}
\hllabelcolour{HlM}
\caption{\hltint{HlM}{Proportional-odds tests and selection under multiplicity corrections.}}
\label{tab:a_ordinal}
\footnotesize
\renewcommand{\arraystretch}{1.15}
\begin{tabularx}{\textwidth}{Xrrrl}
\toprule
\headerrow
\hcell{Test} & \hcell{$\chi^2$} & \hcell{d.f.} & \hcell{$p$} & \hcell{Conclusion} \\
\midrule
\addlinespace
\rowcolor{PanelBg}
\multicolumn{5}{l}{\textbf{Panel A. Tests of proportional odds}} \\
\addlinespace[2pt]
Omnibus likelihood ratio & 74.68 & 35 & $<0.001$ & Rejects \\
\altrow
Wald form, for comparison & 30.54 & 35 & 0.683 & Does not reject \\
\addlinespace
\rowcolor{PanelBg}
\multicolumn{5}{l}{\textbf{Panel B. Selection under each multiplicity correction}} \\
\addlinespace[2pt]
\midrule
\headerrow
\hcell{Correction} & \hcell{Coefficients freed} & \hcell{$\Delta$BIC} & \hcell{} & \hcell{Selected model} \\
\midrule
None & 12 & 14.28 & & Ordered logit \\
\altrow
Benjamini and Hochberg & 9 & 8.83 & & Ordered logit \\
Holm & 2 & $-$14.64 & & Partial proportional odds \\
\bottomrule
\end{tabularx}
\begin{tablenotes}[flushleft]\footnotesize
\item \textit{Note:} A negative $\Delta$BIC favors the partial proportional-odds model over the ordered logit.
\end{tablenotes}
\end{threeparttable}
\end{table}
